\documentclass[10pt,a4paper]{amsart}
\usepackage[margin=19mm]{geometry}
\usepackage{amsmath,amssymb,mathtools}
\usepackage[scr=rsfs,cal=euler]{mathalfa}
\usepackage{xfrac}
\usepackage[unicode,hidelinks,bookmarks=false]{hyperref}
\newcommand{\ie}{i.e.\ }
\newcommand{\cf}{cf.\ }
\newcommand{\resp}{respectively\ }
\newcommand{\Subsection}{\subsection}
\providecommand{\nobreakdash}{\nobreak\hbox{-}}
\setlength{\emergencystretch}{2em}
\multlinegap0pt
\usepackage{bm}
\usepackage{bbm}
\usepackage{dsfont}
\usepackage{xspace}
\usepackage{cancel}
\usepackage{mathtools}
\usepackage{amsthm}
\makeatletter
\@ifundefined{theorem}{\newtheorem{theorem}{Theorem}}{}
\@ifundefined{lemma}{\newtheorem{lemma}[theorem]{Lemma}}{}
\@ifundefined{proposition}{\newtheorem{proposition}[theorem]{Proposition}}{}
\makeatother
\newcommand{\NN}{\mathbb{N}}
\newcommand{\calA}{\mathcal{A}}
\newcommand{\calB}{\mathcal{B}}
\newcommand{\calM}{\mathcal{M}}
\newcommand{\calN}{\mathcal{N}}
\newcommand{\calO}{\mathcal{O}}
\newcommand{\norm}[1]{\lvert {#1} \rvert}
\newcommand{\scrA}{\mathscr{A}}
\newcommand{\scrB}{\mathscr{B}}
\newcommand{\scrN}{\mathscr{N}}
\newcommand{\scrX}{\mathscr{X}}
\newcommand{\sprod}[2]{\langle {#1}, {#2} \rangle}
\title[Mass Functions and Asymptotic Behavior of Caloric Functions ]{Mass Functions and Asymptotic Behavior of Caloric Functions on Affine Buildings}
\author{Effie Papageorgiou, Bartosz Trojan}
\date{Source: arXiv:2506.17042v1 (CC BY 4.0)}
\begin{document}
\maketitle

\noindent\emph{Source and licence.} Mass Functions and Asymptotic Behavior of Caloric Functions on Affine Buildings. Version arXiv:2506.17042v1 (CC BY 4.0), distributed under the Creative Commons Attribution 4.0 International licence, as recorded in the arXiv licence metadata for this exact version and shown on the abstract page https://arxiv.org/abs/2506.17042v1. Full rights notice: licenses/arxiv-2506.17042-CC-BY-4.0.txt.\par\smallskip

\noindent\textit{Editorial scope.} Counted unit: the Theorem whose source label is \texttt{thm:2} in the article (source0.tex lines 1483--1498, source label \texttt{thm:2}), delivered together with the article's own complete original proof of it (source0.tex lines 1499--1633, one \texttt{proof} environment of 135 lines). No mathematics is written, repaired or re-derived: statement and proof are the article's own text line for line. Required context retained uncounted: prop:1 (lines 971--985); eq:33b (lines 1079--1082); eq:3b (lines 1115--1121); lem:4 (lines 1180--1191); lem:1 (lines 1209--1232); eq:38 (lines 1464--1469). Every definition, display and label of those lines is retained unchanged. Omitted from this package and not redistributed: 1--970, 986--1078, 1083--1114, 1122--1179, 1192--1208, 1233--1463, 1470--1482, 1634--3173. Nothing inside a retained excerpt is omitted or altered: every retained body is delivered exactly as the article's lines are written, so no omission receipt of any kind accompanies this package. Citations: the retained excerpts carry no citation command at all, so no bibliography entry of the article is transcribed and no reference is redistributed. Reference adaptation: none was needed and none was made; every reference inside the delivered unit resolves inside it, so no unresolved reference or citation is produced. Printed slips kept literal: no printed word, symbol, hypothesis or number of the counted statement or of its proof was altered. Layout and class: the article's own class is \texttt{amsart} with packages including mathrsfs, amsfonts, amssymb, amsxtra, dsfont, color, xcolor, graphicx, enumitem, babel, fontenc, newtxmath, newtxtext, csquotes; the wrapper is the lane's pinned \texttt{amsart} editorial wrapper at 10pt on A4, which restates the theorem-like environments the retained text needs and adds the article's own macro definitions that the retained text needs (\texttt{\textbackslash NN}, \texttt{\textbackslash calA}, \texttt{\textbackslash calB}, \texttt{\textbackslash calM}, \texttt{\textbackslash calN}, \texttt{\textbackslash calO}, \texttt{\textbackslash norm}, \texttt{\textbackslash scrA}, \texttt{\textbackslash scrB}, \texttt{\textbackslash scrN}, \texttt{\textbackslash scrX}, \texttt{\textbackslash sprod}), with extra wrapper packages (bm, bbm, dsfont, xspace, cancel, mathtools). The delivered numbering is the unit's own. Assets: no retained span contains a figure environment, a graphics-inclusion command, a PDF-page inclusion, a table or a TikZ picture; no image file is copied into this package, the package declares no asset dependency and no image file is redistributed with it. Licence: version arXiv:2506.17042v1 of this article is distributed under the Creative Commons Attribution 4.0 International licence, as recorded in the arXiv licence metadata for this exact version and shown on the abstract page https://arxiv.org/abs/2506.17042v1. A full rights notice is delivered at licenses/arxiv-2506.17042-CC-BY-4.0.txt. Characters: every retained excerpt is pure ASCII, so the delivered engine, XeLaTeX with its default Latin Modern text font, reports no missing character.
\begin{proposition} 
	\label{prop:1}
	We have
	\[
		\mathbf{\Phi}(\lambda)
		\approx \chi_0(\lambda)^{-\frac{1}{2}} \prod_{\alpha \in \Phi^{++}}(1+\langle\alpha, \lambda\rangle).
	\]
	Furthermore,
	\[
		\mathbf{\Phi}(\lambda)
		= \text{ const. } \chi_0(\lambda)^{-\frac{1}{2}}  
		\prod_{\alpha \in \Phi^{++}} \langle\alpha, \lambda\rangle + o(1)
	\]
	as $\langle\alpha, \lambda\rangle \rightarrow +\infty$  for all  $\alpha \in \Phi^{+}.$
\end{proposition}

\begin{equation}
	\label{eq:33b}
	\phi(\delta) \approx \norm{\delta}^2, \qquad \delta \in \calM.
\end{equation}

\begin{equation}
	\label{eq:3b}
	k(n; y, x_n) =
	n^{-\frac{r}{2}-\left|\Phi^{++}\right|} \varrho^n e^{-n \phi\left(\delta_n\right)} 
	\mathbf{\Phi}(\sigma(y, x_n))
	\left(C_0+\mathcal{O}(\left|\delta_n\right|)+\mathcal{O}(n^{-1})\right).
\end{equation}

\begin{lemma}
	\label{lem:4}
	Let $k$ be a transition kernel of an admissible random walk on good vertices of an affine building $\scrX$.
	Let $p \in [1, \infty)$. For all $b > a > 0$, there is $C > 0$ such that for each $n \in \NN$,
	\[
		\sum_{x \in V_g : a\leqslant |\sigma(o,x)|\leqslant b} k(n; x)^p
		\leqslant
		C
		\varrho^{np} \sum_{\lambda \in P^+ : a\leqslant |\lambda|\leqslant b}
		\exp\Big\{p n (\sprod{s_p}{\lambda/n} - \phi(\lambda/n))\Big\}.
	\]
\end{lemma}

\begin{lemma}
	\label{lem:1}
	Let $c > 0$, $\gamma > 0$, $\beta \in \NN^r$ and $x_0 \in \mathfrak{a}$. For each $n \in \NN_0$ we set
	\[
		\Lambda_n = \Big\{\lambda \in P^+ : \norm{\lambda - n x_0} \leqslant n^{\frac{1}{2}+\gamma}\Big\}.
	\]
	Then
	\begin{equation}
		\label{eq:23}
		\sum_{\lambda \in \Lambda_n}
		\Big(\prod_{j = 1}^r |\sprod{\lambda - n x_0}{\alpha_j}|^{\beta_j} \Big)
		e^{-c \frac{\norm{\lambda - n x_0}^2}{n} }
		\approx n^{\frac{r+|\beta|}{2}},
	\end{equation}
	and
	\begin{equation}
		\label{eq:24}
		\sum_{\stackrel{\lambda \in \Lambda_n}{|\lambda - n x_0|_\infty \geqslant n^{1/2}}}
		\Big(\prod_{j = 1}^r |\sprod{\lambda - n x_0}{\alpha_j}|^{\beta_j} \Big)
		e^{-c \frac{\norm{\lambda - n x_0}^2}{n} }
		\approx n^{\frac{r+|\beta|}{2}}
	\end{equation}
	as $n$ tends to infinity uniformly with respect to $x_0 \in \mathfrak{a}$.
\end{lemma}

\begin{equation}
	\label{eq:38}
	0 < \gamma < \frac{1}{4|\Phi^{++}|}
	\quad\text{and}\quad
	\gamma' = 2 \gamma |\Phi^{++}|.
\end{equation} 

\begin{theorem}
	\label{thm:2}
	Let $k_n$ be a transition kernel of an admissible random walk on good vertices of an affine building $\scrX$.
	Then
	\[
		\|k_n \|_{\ell^2}
		\approx 
		n^{-\frac{r}{4}-\frac{|\Phi^{++}|}{2}} \varrho^{n}.
	\]
	Furthermore,
	\[
		\frac{1}{\| k_n\|_{\ell^2}}
		\bigg( \sum_{x \in V_g \setminus \scrN_n^2} k(n; x)^2 \bigg)^\frac{1}{2} = 
		\calO\Big(n^{-\gamma|\Phi^{++}|}\Big).
	\]
\end{theorem}

\begin{proof} 
	We first determine the asymptotic behavior of the norm restricted to $\calN_n^2$. For $v \in \mathfrak{a}$ we set
	\[
		m(v) = \min \big\{\sprod{v}{\alpha} : \alpha  \in \Phi^{++}\big\}.
	\]
	For $n \in \NN$, let
	\begin{align*}
		\calA_n &= \Big\{\lambda \in P^+ : n^{\frac{1}{2}-\gamma} \leqslant |\lambda| \leqslant n^{\frac12+\gamma} \Big\},
		&\calB_n &= \Big\{\lambda \in P^+: |\lambda| \leqslant n^{\frac12+\gamma} \Big\}, \\
		\scrA_n &= \{x \in V_g : \sigma(o, x) \in \calA_n\}, 
		&\scrB_n &= \{x \in V_g : \sigma(o, x) \in \calB_n\}. 
	\end{align*}
	Applying the asymptotic formula \eqref{eq:3b} together with Proposition \ref{prop:1}, we get
	\begin{align*}
		\sum_{x \in \scrN_n^2} k(n; x)^2 
		&\approx
		\varrho^{2n} n^{-r - 2 |\Phi^{++}|} 
		\sum_{x \in \scrN_n^2} e^{- 2 n \phi(\delta)} \chi_0(\sigma(o, x))^{-1} 
		\prod_{\alpha \in \Phi^{++}} (1 + \sprod{\alpha}{\sigma(o, x)}^2)  \\
		&\approx
		\varrho^{2n} n^{-r - 2 |\Phi^{++}|} 
		\sum_{\lambda \in \calN_n^2} 
		e^{-2 n \phi(\delta)} \prod_{\alpha \in \Phi^{++}} (1 + \sprod{\alpha}{\lambda}^2).
	\end{align*}
	Now, in view of \eqref{eq:33b}, there are $c' \geqslant c > 0$ such that 
	\begin{align*}
		\sum_{\lambda \in \calN_n^2} 
		e^{-2 n \phi(\delta)} \prod_{\alpha \in \Phi^{++}} 
		(1 + \sprod{\alpha}{\lambda}^2)
		\leqslant
		\sum_{\lambda \in \calB_n} e^{-c \frac{|\lambda|^2}{n}}
		\prod_{\alpha \in \Phi^{++}}
		(1 + \sprod{\alpha}{\lambda}^2),
	\end{align*}
	and
	\begin{align*}
		\sum_{\lambda \in \calN_n^2} 
		e^{-2 n \phi(\delta)} \prod_{\alpha \in \Phi^{++}} (1 + \sprod{\alpha}{\lambda}^2)
		&\geqslant
		\sum_{\lambda \in \calA_n : m(\lambda) > 1 } e^{-c' \frac{|\lambda|^2}{n}}
		\prod_{\alpha \in \Phi^{++}}
		(1 + \sprod{\alpha}{\lambda}^2).
	\end{align*}
	Observe that if $\lambda \in \calA_n$ and $m(\lambda) > 1$, then $\lambda = \lambda' + \rho$ and 
	\[
		\frac{1}{2} n^{\frac12-\gamma} \leqslant |\lambda'| \leqslant 2 n^{\frac12+\gamma}.
	\]
	Therefore, by Lemma \ref{lem:1}, we get
	\[
		\sum_{\lambda \in \calN_n^2} 
		e^{-2 n \phi(\delta)} \prod_{\alpha \in \Phi^{++}} (1 + \sprod{\alpha}{\lambda})^2
		\approx
		n^{\frac{r}{2} + |\Phi^{++}|},
	\]
	and consequently,
	\begin{equation}
		\label{eq:34}
		\sum_{x \in \scrN_n^2} k(n; x)^2 \approx \varrho^{2n} n^{-\frac{r}{2} - |\Phi^{++}|}.
	\end{equation}
	Next, let us consider the region $V_g \setminus \scrB_n$. By Lemma \ref{lem:4} and \eqref{eq:33b}, there is $c > 0$ 
	such that 
	\begin{align*}
		\sum_{x \in V_g \setminus \scrB_n}
		k(n; x)^2 
		&\lesssim 
		\varrho^{2n} 
		\sum_{\lambda \in P^+ \setminus \calB_n} e^{-c \frac{|\lambda|^2}{n}} \\
		&\lesssim
		\varrho^{2n}
		e^{-\frac{c}{2} n^{2\gamma}} \sum_{\lambda \in P} e^{-c \frac{|\lambda|^2}{2n}},
	\end{align*}
	thus
	\begin{equation}
		\label{eq:36}
		\sum_{x \in V_g \setminus \scrB_n} k(n; x)^2
		\lesssim \varrho^{2n} n^{-\frac{r}{2} - |\Phi^{++}|} e^{-\frac{c}{4} n^{2\gamma}}.
	\end{equation}
	Let us observe that in the region $x \in \scrB_n \setminus \scrA_n$, for each $\alpha \in \Phi^{++}$, 
	we have
	\[
		1 + \sprod{\lambda}{\alpha}^2 \lesssim n^{1- 2 \gamma}, 
	\]
	hence by the asymptotic formula \eqref{eq:3b} together with Proposition \ref{prop:1} and \eqref{eq:33b}, we get
	\begin{align*}
		\sum_{x \in \scrB_n \setminus \scrA_n} k(n; x)^2
		&\lesssim
		\varrho^{2n} n^{-r - 2 |\Phi^{++}|} 
		\sum_{x \in \scrB_n \setminus \scrA_n}
		e^{-2n \phi(\delta)} \chi_0(\sigma(o, x))^{-1} 
		\prod_{\alpha \in \Phi^{++}} \big(1 + \sprod{\sigma(o, x)}{\alpha}^2\big) \\
		&\lesssim
		\varrho^{2n} n^{-r -|\Phi^{++}|-2 \gamma|\Phi^{++}|}
		\sum_{\lambda \in \calB_n \setminus \calA_n} e^{-c \frac{|\lambda|^2}{n}}.
	\end{align*}
	Therefore, we obtain
	\begin{equation}
		\label{eq:37}
		\sum_{x \in \scrB_n \setminus \scrA_n} k(n; x)^2
		\lesssim
		\varrho^{2n} n^{-\frac{r}{2} -|\Phi^{++}|} n^{-2 \gamma|\Phi^{++}|}.
	\end{equation}
	Lastly, we need to consider the region $\scrA_n \setminus \scrN_n^2$. For $x \in \scrA_n \setminus \scrN_n^2$, there
	is $\beta \in \Phi^{++}$ such that
	\[
		\sprod{\sigma(o, x)}{\beta} \leqslant n^{\frac{1}{2}-\gamma'}.
	\]
	However, for all $\alpha \neq \beta$, we have
	\[
		\sprod{\sigma(o, x)}{\alpha} \leqslant C n^{\frac12+\gamma},
	\]
	thus by \eqref{eq:38}
	\[
		\prod_{\alpha \in \Phi^{++}} \big(1 + \sprod{\sigma(o, x)}{\alpha}^2\big)
		\lesssim 
		n^{|\Phi^{++}|} n^{-2(\gamma' - \gamma (|\Phi^{++}| - 1))}
		\leqslant
		n^{|\Phi^{++}|} n^{-2\gamma (|\Phi^{++}| + 1)}.
	\]
	Therefore, by the asymptotic formula \eqref{eq:3b} together with Proposition \ref{prop:1} and \eqref{eq:33b}, we obtain
	\begin{align}
		\nonumber
		\sum_{x \in \scrA_n \setminus \scrN_n^2}
		k(n; x)^2
		&\lesssim
		\varrho^{2n} n^{-r - 2 |\Phi^{++}|} 
		\sum_{\lambda \in \calA_n \setminus \calN_n^2}
		e^{-2n \phi(\delta)}
		\prod_{\alpha \in \Phi^{++}} \big(1 + \sprod{\lambda}{\alpha}^2\big) \\
		\label{eq:39}
		&\lesssim
		\varrho^{2n} n^{-\frac{r}2 - |\Phi^{++}|}  n^{-2\gamma (|\Phi^{++}| + 1)}.
	\end{align}
	Using estimates \eqref{eq:36}, \eqref{eq:37} and \eqref{eq:39} together with the asymptotic \eqref{eq:34} 
	we complete the proof.
\end{proof}

\end{document}
